\bookchapter{第四章：安裝與管理App——手機之“技能包”}

\chapterintro{App使手機無所不能，然亦須管之。猶汝予新員工安排工位、分配權限也，裝App亦須一套規範之流程。}

\booksection{一、官方之店：至安之“應用超市”}

\begin{itemize}
  \item \textbf{安卓}：Google Play（國外）或各品牌自帶之店（如華爲應用市場、小米應用商店）。
  \item \textbf{iOS}：惟 App Store。
  \item \textbf{鴻蒙}：華爲應用市場。
\end{itemize}

核心之安之義：經審覈、更有所保、支付之安。行之指：\textbf{首選惟用官方之店下載！}

\booksection{二、純淨/防護之模：開手機之“防毒牆”}

主流之國產品牌（華爲、小米、OPPO、vivo、榮耀）皆內置了\textbf{純淨之模}或\textbf{防護之模}，此乃系統級別之安之開關。\textbf{強烈建議諸用戶皆開之}。

\booksubsection{其究何物？}

純淨之模（亦稱“防護之模”）之\textbf{核心之功乃阻第三方App之安}：當其開時，系統將\textbf{惟允自官方應用之店裝App}，自瀏覽器、二維碼、聊天軟件等渠道所下之安裝包將直被攔，從源頭杜惡意之軟件。

\booksubsection{適用之人}

\begin{itemize}
  \item \textbf{年老之用戶}：分不清"官方"與"假冒"之告，純淨之模乃至穩之護符。
  \item \textbf{對應用之權不敏之用戶}：不需裝小衆App、測試版、外服遊戲。
  \item \textbf{欲升設備安之諸用戶}：縱爲發燒友，開之亦勝於不——真有臨需，再獨閉可也。
\end{itemize}

\booksubsection{開啓之徑（以華爲/小米爲例）}

\begin{itemize}
  \item \textbf{華爲/鴻蒙}：[設置] > [系統和更新] >\textbf{[純淨之模]}，擇"開"。
  \item \textbf{小米}：[設置] > [隱私保護] >\textbf{[保護隱私]}或 [設置] > [應用設置] > [應用安裝]，閉"安裝未知應用"。
  \item \textbf{OPPO/vivo}：[設置] > [安全] >\textbf{[安裝外部來源應用]}，逐個閉授權。
\end{itemize}

\booksubsection{應用之店之安檢測之機}

應用之店非“擺設”，其背有完之安之流水線：

\begin{itemize}
  \item \textbf{病毒之掃}：每App上架前經數家殺毒引擎聯合掃，識已知木馬與病毒之特徵。
  \item \textbf{權限之核}：店將審App所請之權限是否"合"——一電筒App要通訊錄之權？直駁或下架。
  \item \textbf{開發者之認}：上架須企業之資、個人身份證或信用卡之認。出問題能尋人，非"黑App"。
  \item \textbf{持續之監}：已上架之App亦被持續復掃，新發之惡意行將強下架。
\end{itemize}

\booksection{三、第三方之來與“電子寄生蟲”：險之極高！}

自網頁或不明之鏈下App，或夾帶病毒、竊隱私、盜刷支付。此處吾等特提一新概：\textbf{“電子寄生蟲”}。

\textbf{何爲“電子寄生蟲”？}\par

指彼等\textbf{被植木馬之“破解版”軟件}——汝以爲“省了錢”，實則“中了招”。常見之形：

\begin{itemize}
  \item "破解版"、"免費版"、"去告版"付費軟件；
  \item "破解版"、"免費版"、"去告版"付費軟件；
  \item 羣聊中"小告"所薦之所謂"內部版"App。
\end{itemize}

此等程序表面能用，背或在\textbf{偷通訊錄、讀短信、監麥克風、後臺挖礦}。純淨之模乃防之\textbf{第一道防線}——直不使汝裝。

\booksubsection{第二道防線之建議：}

\begin{itemize}
  \item \textbf{定期全盤病毒掃}：用手機自帶之安中心（如華爲手機管家、小米安中心）每月掃一。
  \item \textbf{慎予敏感之權}：通訊錄、短信、位置、相機、麥克風——此五"高危之權"必嚴審。
  \item \textbf{關注耗電與流量之異}：手機突燙、電速掉、流量半夜跑——八成後臺藏"寄生蟲"。
\end{itemize}

若萬不得已須裝第三方App：

\begin{enumerate}
  \item 保來源絕對可信。
  \item 保來源絕對可信。
  \item 事後用安App掃。
\end{enumerate}

終極之建議：於系統設置中，閉“安裝未知應用”之權。

\booksection{四、裝、更、卸}

\begin{itemize}
  \item \textbf{裝}：店中點"獲取"或"裝"，裝畢首開將問汝要權。
  \item \textbf{更}：\textbf{必更！}主爲修安之洞。建議開自動更。
  \item \textbf{卸}：長按圖標拖至卸區。欲清更徹，往 [設置] > [應用管理] 中尋而卸之。
\end{itemize}

\booksection{五、權限之管：汝乃App之“老闆”}

App要權（如通訊錄、位置）有時合（地圖要位置），有時不合（電筒要通訊錄）。

\begin{itemize}
  \item \textbf{權限之擇}："始終允"、"僅用時允"、"禁止"、"每問"。
  \item \textbf{管之則}：
  \begin{enumerate}
    \item 最小化：非必不與。
    \item 最小化：非必不與。
    \item 情境化：似打車軟件，惟用時予位置之權。
  \end{enumerate}
\end{itemize}

\begin{tcbwarn}{重點盯防之權}
\begin{itemize}
  \item 通訊錄/短信：或被用以發垃圾短信。
  \item 實時位置：或被追蹤行蹤。
  \item 實時位置：或被追蹤行蹤。
  \item 文件訪問：或讀汝隱私之文件。
\end{itemize}

\end{tcbwarn}

\booksection{六、軟件安裝包之構成要素}

汝曾否思之，自應用之店所下之安裝包，內究裝何物？知安裝包之構成，可助汝更解App之作，亦能於遇問題時知何排查。

\booksubsection{安卓APK文件之結構}

APK（Android Package Kit）本乃一ZIP壓縮包，解壓可見以下核心之目錄與文件：

\begin{itemize}
  \item \textbf{AndroidManifest.xml}：應用之"身份證"與"說明書"。錄應用之包名、版本號、所需之權、支持之系統版、入口Activity等核心之信。系統裝應用時首讀此文件。
  \item \textbf{classes.dex}：應用之核心碼。Java/Kotlin源代碼編譯後變Dalvik字節碼，存此文件中。一應用或有多dex文件（如classes2.dex、classes3.dex）。
  \item \textbf{resources.arsc}：編譯後之資源索引表。錄所有資源之ID與對應之關係。
  \item \textbf{res/目錄}：資源之文件夾。含圖、佈局文件、字符串、顏色定義等。不同辨率之圖放不同drawable文件夾中（如drawable-hdpi、drawable-xhdpi）。
  \item \textbf{assets/目錄}：原始資源之文件夾。存不編譯之資源，如網頁文件、字體、音頻等。App運行時可經AssetManager讀之。
  \item \textbf{lib/目錄}：本地庫之文件夾。存編譯好之.so文件（C/C++碼編譯之產），按CPU架構分類（如arm64-v8a、armeabi-v7a）。遊戲與需高性能算之App通常有此目錄。
  \item \textbf{META-INF/目錄}：簽名信之文件夾。含MANIFEST.MF、CERT.SF、CERT.RSA三文件，用於驗APK之完與簽名。
\end{itemize}

\booksubsection{iOS IPA文件之結構}

IPA（iOS App Store Package）同爲壓縮包，結構略異：

\begin{itemize}
  \item \textbf{Payload/目錄}：核心內容之所在。內有.app文件夾，含應用所有內容。
  \item \textbf{Info.plist}：應用之配置文件，似安卓之AndroidManifest.xml。錄應用ID、版本、權之需、支持之設備等。
  \item \textbf{可執行文件}：與.app文件夾同名之主程序，乃編譯後之Mach-O格式二進制文件。
  \item \textbf{Frameworks/目錄}：應用所賴之動態庫與框架。
  \item \textbf{資源文件}：圖、音、視頻等資源，通常存於.app文件夾內。
  \item \textbf{\_CodeSignature/目錄}：碼之簽名信，用於驗應用之完與源。
\end{itemize}

\booksubsection{簽名驗證之機}

何故自官方之店所下之App更安？簽名驗證爲關鍵。

\begin{itemize}
  \item \textbf{安卓之籤}：開發者以私鑰對APK籤，系統裝時以公鑰驗。若APK被篡（如植病毒），簽名即失效，系統拒裝。Google Play猶須App Signing，由Google保管簽名之鑰，進一步護應用之安。
  \item \textbf{iOS之籤}：蘋果之審更嚴。每IPA必具蘋果所頒之證書籤，否則不能裝。此即iOS惟自App Store或企業證書裝應用之由。企業證書籤之應用有效期有限，且蘋果可隨吊銷。
  \item \textbf{籤之作}：
  \begin{enumerate}
    \item 驗源：認應用確來自所聲之開發者。
    \item 驗完：保應用於傳中未被篡。
    \item 驗完：保應用於傳中未被篡。
  \end{enumerate}
\end{itemize}

\booksubsection{安裝包之大與優化}

汝或發見，應用之店所顯之App大與裝後佔之空間不一。蓋因：

\begin{itemize}
  \item \textbf{下載之大}：壓縮後之APK/IPA大。
  \item \textbf{安裝之大}：解壓後之大，通常遠大於下載之大。
  \item \textbf{運行時之佔}：App運行時猶生緩存、數據文件，佔更多之空間。
\end{itemize}

現代App之優化之方：

\begin{itemize}
  \item \textbf{App Bundle（安卓）}：Google Play之動態分發之技。用戶下載時惟得適己設備之資源（如惟下載對應屏辨率之圖），大減下載之大。
  \item \textbf{App Thinning（iOS）}：類之技，蘋果將據設備自動裁不需之資源。
  \item \textbf{按需之下}：遊戲等大型應用可先下核心包，後續再下資源包。
\end{itemize}

\begin{tcbinfo}{小知識：何有之App安裝包特大？}
小知識：何有之App安裝包特大？

\end{tcbinfo}

\begin{tcbpractice}{實操之環：審三App之權}
實操之環：審三App之權

\end{tcbpractice}
